\documentclass[11pt]{article}
\usepackage[letterpaper,margin=0.85in]{geometry}\usepackage[T1]{fontenc}\usepackage[utf8]{inputenc}\usepackage{amsmath,amssymb,enumitem,booktabs,array,xcolor,fancyhdr,tabularx}
\setlist{itemsep=2pt,topsep=3pt}\pagestyle{fancy}\fancyhf{}\renewcommand{\headrulewidth}{0.3pt}
\fancyhead[L]{\small 2.8 The Quotient Rule}\fancyhead[R]{\small Make-up key}\fancyfoot[C]{\small\thepage}
\definecolor{ans}{RGB}{28,125,74}\newcommand{\A}[1]{\textcolor{ans}{\textbf{#1}}}
\begin{document}

\section*{2.8 The Quotient Rule}
\textbf{You missed this lesson. Copy this into your notes. That is all you need to do.}\par\medskip
\subsection*{Question of the day}
How do you differentiate one function divided by another?\par
\subsection*{Targets}\begin{itemize}\item I can state the quotient rule.
\item I can differentiate a quotient and find its slope at a point.
\item I can use a table of values with the rule.\end{itemize}
\subsection*{Concepts}\begin{itemize}
\item \textbf{The quotient rule} (h? = (g f? -- f g?) / g\textsuperscript{2}): Bottom times top?, minus top times bottom?. All over the bottom squared.
\item \textbf{A false rule} (h? ? f? / g?): The derivative of a quotient is not f?/g?. Check it against the tangent slope.
\item \textbf{Choosing a method} ((3x\textsuperscript{2} + 2?x)/x = 3x + 2x\textsuperscript{-1}?\textsuperscript{2}): Simplify first when the bottom is one term. Check with a calculator derivative.
\end{itemize}
\subsection*{What the class predicted, and what is true}\begin{itemize}
\item A student guesses that for $h=\dfrac{\sin x}{x^{2}}$ the derivative is $\dfrac{\cos x}{2x}$. At $x=1$ that gives $+0.27$, but the graph of $h$ is falling there. The guess is... \A{wrong: the sign disagrees with the graph.} The graph falls at x = 1, so the slope is negative. A positive value cannot be right.
\item The curve $y=\dfrac{1}{x}$ for $x>0$ falls as $x$ increases. Its slope at $x=2$ is... \A{negative.} The curve is falling, so the slope is negative.
\end{itemize}
\section*{Huddle notes (answers filled in)}
\subsection*{Definitions}\begin{itemize}
\item \textbf{Quotient rule:} If h = f/g, then h? = (\A{g f?} -- \A{f g?}) / \A{g\textsuperscript{2}}.
\item \textbf{Order:} The \A{bottom} times the derivative of the \A{top} comes first.
\item \textbf{Shortcut:} If the bottom is one term, \A{simplify} first and use the power rule.
\end{itemize}
\subsection*{Worked example}
Find h?(5) for h(x) = (x\textsuperscript{2} + 1)/(x -- 3).\\
\begin{itemize}
\item g f? -- f g? = (x -- 3)(2x) -- (x\textsuperscript{2} + 1)(1)
\item At x = 5: 2 $\cdot$ 10 -- \A{26} = \A{--6}
\item Divide by g\textsuperscript{2} = \A{4}: h?(5) = \A{--3/2}
\end{itemize}
\subsection*{Label the rule}
\begin{enumerate}
\item The numerator is \A{bottom times top?} minus top times bottom?.
\item Order matters: \A{subtraction} is not commutative.
\item The denominator is the bottom \A{squared}.
\end{enumerate}
\section*{One worked sample of each skill}
\subsection*{Differentiate a quotient}
\noindent\textbf{Level 1.} Differentiate with the quotient rule.\\ \A{Answer: $\frac{-5}{\left(x - 2\right)^{2}}$} \quad Numerator: $(x-2)(1)-(x+3)(1)=-5$. So $\dfrac{-5}{(x-2)^{2}}$.\par\medskip
\noindent\textbf{Level 2.} Differentiate with the quotient rule.\\ \A{Answer: $\frac{x^{2} + 2 x}{\left(x + 1\right)^{2}}$} \quad $\dfrac{(x+1)(2x)-x^{2}}{(x+1)^{2}}=\dfrac{x^{2}+2x}{(x+1)^{2}}$.\par\medskip
\noindent\textbf{Level 3.} Differentiate. Name f and g first.\\ \A{Answer: $\frac{x^{2} \cdot e^{x} - 2 x^{1} \cdot e^{x}}{x^{4}}$} \quad $\dfrac{x^{2}e^{x}-e^{x}\cdot2x^{1}}{x^{4}}$.\par\medskip
\noindent\textbf{Level 4.} Differentiate. You do not have to simplify.\\ \A{Answer: $\frac{x^{5} \cos\left(x\right) - 5 x^{4} \sin\left(x\right)}{x^{10}}$} \quad $\dfrac{x^{5}\cos x-5x^{4}\sin x}{x^{10}}$.\par\medskip
\subsection*{Slope at a point}
\noindent\textbf{Level 1.} Let $h(x)=\dfrac{x}{x+2}$. Find $h'(2)$.\\ \A{Answer: 0.125} \quad $h'(x)=\dfrac{(x+2)-x}{(x+2)^{2}}=\dfrac{2}{(x+2)^{2}}$, so $h'(2)=\dfrac{2}{16}=0.125$.\par\medskip
\noindent\textbf{Level 2.} Let $h(x)=\dfrac{x^{2}+5}{x-1}$. Find $h'(3)$.\\ \A{Answer: -0.5} \quad $h'(3)=\dfrac{(2)(6)-(14)(1)}{4}=-0.5$.\par\medskip
\noindent\textbf{Level 3.} A table gives $f(4)=-5$, $f'(4)=2$, $g(4)=10$ and $g'(4)=-6$. Let $h=\dfrac{f}{g}$. Find $h'(4)$.\\ \A{Answer: -0.1} \quad $h'(4)=\dfrac{(10)(2)-(-5)(-6)}{(10)^{2}}=-0.1$.\par\medskip
\noindent\textbf{Level 4.} A table gives $f(3)=8$, $f'(3)=-5$, $g(3)=10$ and $g'(3)=4$. Let $h=\dfrac{f}{g}$. Find $h'(3)$.\\ \A{Answer: -0.82} \quad $h'(3)=\dfrac{(10)(-5)-(8)(4)}{(10)^{2}}=-0.82$.\par\medskip
\subsection*{Simplify first}
\noindent\textbf{Level 1.} Divide first, then differentiate.\\ \A{Answer: $4$} \quad $\dfrac{4x^{2}+4x}{x}=4x+4$, whose derivative is $4$.\par\medskip
\noindent\textbf{Level 2.} Divide first, then differentiate.\\ \A{Answer: $4 x$} \quad $2x^{2}+7$ has derivative $4x$.\par\medskip
\noindent\textbf{Level 3.} Divide first, then differentiate.\\ \A{Answer: $5 - \frac{2}{2 x^{\frac{3}{2}}}$} \quad $5x+2x^{-1/2}$ has derivative $5-\tfrac{2}{2}x^{-3/2}$.\par\medskip
\noindent\textbf{Level 4.} Divide first, then differentiate.\\ \A{Answer: $5 + \frac{2}{x^{3}}$} \quad $5x+5-1x^{-2}$ has derivative $5+\dfrac{2}{x^{3}}$.\par\medskip
\subsection*{Horizontal tangents}
\noindent\textbf{Level 1.} Let $h(x)=\dfrac{x}{x^{2}+4}$ for $x>0$. At which $x$ is the tangent horizontal?\\ \A{Answer: 2} \quad $h'(x)=\dfrac{(x^{2}+4)-x(2x)}{(x^{2}+4)^{2}}=\dfrac{4-x^{2}}{(x^{2}+4)^{2}}$, which is 0 at $x=2$.\par\medskip
\noindent\textbf{Level 2.} Let $h(x)=\dfrac{x}{x^{2}+4}$ for $x>0$. At which $x$ is the tangent horizontal?\\ \A{Answer: 2} \quad The numerator is $4-x^{2}$, which is 0 at $x=2$.\par\medskip
\noindent\textbf{Level 3.} Let $h(x)=\dfrac{x^{2}+25}{x}$ for $x>0$. At which $x$ is the tangent horizontal?\\ \A{Answer: 5} \quad The numerator is $x^{2}-25$, which is 0 at $x=5$.\par\medskip
\noindent\textbf{Level 4.} Let $h(x)=\dfrac{e^{x}}{x^{6}}$ for $x>0$. At which $x$ is the tangent horizontal?\\ \A{Answer: 6} \quad $h'(x)=\dfrac{e^{x}x^{5}(x-6)}{x^{12}}$, which is 0 at $x=6$ (e? is never 0).\par\medskip
\end{document}